% MBDyn (C) is a multibody analysis code.
% http://www.mbdyn.org
%
% Copyright (C) 1996-2026
%
% Pierangelo Masarati	<pierangelo.masarati@polimi.it>
% Paolo Mantegazza	<paolo.mantegazza@polimi.it>
%
% Dipartimento di Ingegneria Aerospaziale - Politecnico di Milano
% via La Masa, 34 - 20156 Milano, Italy
% http://www.aero.polimi.it
%
% Changing this copyright notice is forbidden.
%
% This program is free software; you can redistribute it and/or modify
% it under the terms of the GNU General Public License as published by
% the Free Software Foundation (version 2 of the License).
%
%
% This program is distributed in the hope that it will be useful,
% but WITHOUT ANY WARRANTY; without even the implied warranty of
% MERCHANTABILITY or FITNESS FOR A PARTICULAR PURPOSE.  See the
% GNU General Public License for more details.
%
% You should have received a copy of the GNU General Public License
% along with this program; if not, write to the Free Software
% Foundation, Inc., 59 Temple Place, Suite 330, Boston, MA  02111-1307  USA
%

\emph{Author: Andrea Zanoni}

\bigskip

This module allows MBDyn to act as the \emph{Flight Dynamics Module} node of
a MARSH (Modular Aircraft Research Simulator Handler,
\htmladdnormallink{\texttt{https://marsh-sim.github.io/}}{https://marsh-sim.github.io/})
flight simulator.
MARSH is a modular simulator framework whose nodes (flight model, pilot
controls, visualisation, instruments, motion platform, \ldots) exchange
MAVLink
(\htmladdnormallink{\texttt{https://mavlink.io/}}{https://mavlink.io/})
messages over UDP, routed by a central component called the \emph{MARSH
Manager}.

The element publishes the state of a structural node as MAVLink messages,
and decodes incoming messages, making the received values available to the
rest of the model as private data.
Which messages are published and which are subscribed to is declared in the
input file.
A \kw{HEARTBEAT} message is sent periodically, and incoming \kw{HEARTBEAT}
messages from the Manager are tracked.

\subsubsection{Syntax}

\begin{Verbatim}[commandchars=\\\{\}]
    \bnt{name} ::= \kw{marsh}

    \bnt{module_data} ::=
        [ \kw{manager address} , (\ty{string}) " \bnt{manager_address} " , ]
        [ \kw{manager port} , (\ty{integer}) \bnt{manager_port} , ]
        [ \kw{local port} , (\ty{integer}) \bnt{local_port} , ]
        [ \kw{system id} , (\ty{integer}) \bnt{system_id} , ]
        [ \kw{component id} , (\ty{integer}) \bnt{component_id} , ]
        [ \kw{heartbeat rate} , (\ty{real}) \bnt{heartbeat_rate} , ]
        [ \bnt{message_clause} [ , ... ] ]
        [ , \bnt{output_flag} ]

    \bnt{message_clause} ::=
        { \kw{send} , (\ty{string}) " \bnt{outbound_message} " \bnt{outbound_data}
        | \kw{subscribe} , (\ty{string}) " \bnt{inbound_message} " }
\end{Verbatim}

\paragraph{Element Ordering}
The elements this one refers to, namely any \kw{aircraft instruments}
element used as the source of an outbound message, must be declared
\emph{before} it; and so must this element with respect to any element that
reads its private data.

\paragraph{Input Parameters}
\begin{itemize}
    \item \nt{manager\_address}: IP address or host name of the MARSH
        Manager (optional; defaults to \kw{"127.0.0.1"});
    \item \nt{manager\_port}: UDP port the MARSH Manager listens on
        (optional; defaults to 24400, the MARSH default);
    \item \nt{local\_port}: UDP port the element binds locally (optional;
        defaults to 0, meaning the operating system assigns an ephemeral
        port);
    \item \nt{system\_id}: MAVLink system ID, which in MARSH identifies the
        simulation session; all the nodes taking part in one simulation
        should share it (optional; defaults to 1);
    \item \nt{component\_id}: MAVLink component ID, which in MARSH
        identifies the role of the node (optional; defaults to 26,
        \kw{MARSH\_COMP\_ID\_FLIGHT\_MODEL});
    \item \nt{heartbeat\_rate}: rate, in Hz, at which \kw{HEARTBEAT}
        messages are sent (optional; defaults to 1.0, the rate MARSH
        currently expects); must be positive;
    \item \nt{output\_flag}: the usual output flag of elements.
\end{itemize}

Each \nt{message\_clause} adds one message to those handled by the element;
clauses may be repeated, in any order, and the same message should not be
declared twice.
Note that the Manager only forwards to a component the messages that are
relevant for its component ID, and that only the first component in a
network with a given system ID and component ID has its messages forwarded;
the messages of any further ones are ignored.

\subsubsection{Outbound Messages}

\begin{Verbatim}[commandchars=\\\{\}]
    \bnt{outbound_message} ::= { \kw{SIM_STATE} | \kw{MOTION_CUE_EXTRA} }

    \bnt{outbound_data} ::=
        [ , \bnt{message_source} ]
        [ , \kw{field} , (\ty{string}) " \bnt{field_name} " ,
            (\hty{ScalarValue}) \bnt{field_value} [ , ... ] ]

    \bnt{message_source} ::=
        { \kw{aircraft instruments} , \bnt{aircraft_instruments_label}
        | \kw{node} , \bnt{node_label}
            [ , \kw{position} , (\hty{Vec3}) \bnt{offset} ]
            [ , \kw{orientation} , (\hty{OrientationMatrix}) \bnt{orientation} ] }
\end{Verbatim}

MAVLink expresses the motion in body axes, with $x$ pointing forward, $y$ to
the right and $z$ down; angles are in radians, angular velocities in rad/s
and accelerations in m/s$^2$, and the ``world'' components in a
North-East-Down frame.
Both messages are sent once per converged time step.

\paragraph{Message Source}
The state published by a message may come either from an
\htmlref{\kw{aircraft instruments}}{sec:EL:AERO:AIRCRAFTINSTRUMENTS}
element, or from a structural node.

Referring to an \kw{aircraft instruments} element is the recommended choice,
and the only one that fills in the whole of \kw{SIM\_STATE}: that element
resolves the orientation of the aircraft with respect to its node,
the flight mechanics angles, and the geodetic position, which the module
converts from the radians of MBDyn to the degrees of MAVLink.
Recall that an \kw{aircraft instruments} element requires an
\kw{air properties} element to be defined.

Referring to a structural node directly is meant for the cases in which no
\kw{aircraft instruments} element is needed, for instance when only
\kw{MOTION\_CUE\_EXTRA} is published.
The point and the reference frame the motion refers to may then be moved
with respect to the node by \nt{offset} and \nt{orientation}; the latter
also is the intended means to reconcile the axes of the node with the
convention MAVLink expects.
The geodetic position and the Euler angles are not available in this case,
and are left null unless the corresponding fields are set explicitly.

\begin{itemize}
    \item \kw{SIM\_STATE} reports the attitude, as a quaternion and, when
        available, as Euler angles, the angular velocity, the acceleration,
        the geodetic position and the ground velocity.
        The standard deviation fields are left null.
    \item \kw{MOTION\_CUE\_EXTRA} is a MARSH specific message intended for
        motion cueing; it reports the angular velocity and the acceleration,
        in the surge, sway and heave sense.
\end{itemize}

Be aware that MAVLink defines the acceleration reported by these messages as
the \emph{specific force}, namely the quantity an accelerometer would
measure, which includes the contribution of gravity; a body at rest thus
reports a vector whose magnitude is that of the gravity acceleration, rather
than a null one.
The element accounts for this, subtracting the gravity acceleration of the
\kw{gravity} element, if one is defined, from that of the aircraft.

\paragraph{Fields}
Any field of a message may be set explicitly to \nt{field\_value}, which is
an arbitrary \hty{ScalarValue}, namely either
\kw{drive} , (\hty{DriveCaller}) \nt{drive\_caller}
or
\kw{node dof} , (\hty{ScalarDof}) \nt{scalar\_dof}.
An explicitly set field overrides whatever the message source computed for
it, so this serves both to complete what a source leaves out, and to publish
quantities that the dynamics of the aircraft does not describe, such as
special effects or the simulation of a faulty sensor.
Since a source is optional, a message may as well be assembled out of
\kw{field} clauses alone.

The \nt{field\_name} is that of the corresponding MAVLink message field:
\kw{q1} to \kw{q4}, \kw{roll}, \kw{pitch}, \kw{yaw},
\kw{xacc}, \kw{yacc}, \kw{zacc}, \kw{xgyro}, \kw{ygyro}, \kw{zgyro},
\kw{lat}, \kw{lon}, \kw{alt}, \kw{std\_dev\_horz}, \kw{std\_dev\_vert},
\kw{vn}, \kw{ve} and \kw{vd} for \kw{SIM\_STATE};
\kw{vel\_roll}, \kw{vel\_pitch}, \kw{vel\_yaw},
\kw{acc\_x}, \kw{acc\_y} and \kw{acc\_z} for \kw{MOTION\_CUE\_EXTRA}.
Note that \kw{lat} and \kw{lon} are in degrees, as MAVLink defines them.

\subsubsection{Inbound Messages}

\begin{Verbatim}[commandchars=\\\{\}]
    \bnt{inbound_message} ::= \kw{MANUAL_CONTROL}
\end{Verbatim}

\begin{itemize}
    \item \kw{MANUAL\_CONTROL} carries the pilot control inputs, as produced
        by the MARSH controls node.
        Its axes and buttons are made available as private data.
\end{itemize}

\subsubsection{Output}

This element sends output to the \texttt{.usr} file.
Currently, each entry contains
\begin{itemize}
\item[1)] the label
\item[2)] a flag, 1 if a \kw{HEARTBEAT} message has been received from the
    MARSH Manager, 0 otherwise
\end{itemize}

\subsubsection{Private Data}

Each subscribed message contributes its own private data, named after the
message.
Subscribing to \kw{MANUAL\_CONTROL} makes the following data available:
\begin{enumerate}
\item \kw{"MANUAL\_CONTROL.x"} the $x$ axis of the control input
\item \kw{"MANUAL\_CONTROL.y"} the $y$ axis of the control input
\item \kw{"MANUAL\_CONTROL.z"} the $z$ axis of the control input
\item \kw{"MANUAL\_CONTROL.r"} the $r$ axis of the control input
\item \kw{"MANUAL\_CONTROL.buttons"} the button bitfield
\end{enumerate}
The four axes are those of the MAVLink message, normalized by the sender in
the range $[-1000, 1000]$, and are passed on unscaled; in the MARSH
convention they respectively correspond to roll, pitch, thrust and yaw
control.
They can be accessed with the \kw{element} drive caller, to use the pilot
input elsewhere in the model.

\subsubsection{Notes}

MBDyn is expected to run in real time when used as a node of a simulator;
see the \kw{real time} card in the \kw{initial value} problem block.
Messages are sent and received once per converged time step, so the rate at
which they are exchanged follows the time step of the simulation, and the
periodic \kw{HEARTBEAT} message is timed on the simulation time.

Further details, including the MARSH specific MAVLink dialect the module is
built against and how to extend it with further messages, are documented in
\texttt{modules/module-marsh/README.md}.
